% =====================================================================
%  The Solver Tolerance Is Not a Solver Parameter
%
%  Self-contained. Cites only the standard published literature.
% =====================================================================
\documentclass[11pt,a4paper]{article}

\usepackage[utf8]{inputenc}
\usepackage[T1]{fontenc}
\usepackage{lmodern}
\usepackage[a4paper,margin=2.4cm]{geometry}
\usepackage{amsmath,amssymb,amsthm,mathtools}
\usepackage{bm}
\usepackage{booktabs}
\usepackage{array}
\usepackage{enumitem}
\usepackage{microtype}
\usepackage{xcolor}
\usepackage{listings}
\usepackage{algorithm}
\usepackage{algpseudocode}
\usepackage[round,sort&compress]{natbib}
\usepackage[colorlinks=true,linkcolor=blue!45!black,
            citecolor=blue!45!black,urlcolor=blue!45!black]{hyperref}
\usepackage{cleveref}

% ---- Theorem environments ----
\newtheoremstyle{pstyle}{\topsep}{\topsep}{\itshape}{}{\bfseries}{.}{.5em}{}
\theoremstyle{pstyle}
\newtheorem{theorem}{Theorem}[section]
\newtheorem{lemma}[theorem]{Lemma}
\newtheorem{proposition}[theorem]{Proposition}
\newtheorem{corollary}[theorem]{Corollary}
\theoremstyle{definition}
\newtheorem{definition}[theorem]{Definition}
\newtheorem{assumption}[theorem]{Assumption}
\newtheorem{criterion}[theorem]{Criterion}
\newtheorem{protocol}[theorem]{Protocol}
\newtheorem{rulename}[theorem]{Rule}
\theoremstyle{remark}
\newtheorem{remark}[theorem]{Remark}

\crefname{criterion}{Criterion}{Criteria}
\Crefname{criterion}{Criterion}{Criteria}
\crefname{protocol}{Protocol}{Protocols}
\Crefname{protocol}{Protocol}{Protocols}
\crefname{rulename}{Rule}{Rules}
\Crefname{rulename}{Rule}{Rules}

% ---- listing style ----
\definecolor{kw}{RGB}{30,64,175}
\definecolor{cmt}{RGB}{22,101,52}
\definecolor{strc}{RGB}{159,18,57}
\definecolor{lstbg}{RGB}{250,250,250}
\lstdefinelanguage{cfc}{
  morekeywords={floor,circuit,species,reaction,solve,yield,patient,probe,
    panel,measure,localize,using,until,resolved,otherwise,decline,cell,
    gap,close,toward,coherent,basis,assert,emit,report,let,import,
    holonomy,loops,of,wegscheider,check,from,where,in,foreach,scan,
    against,collect,tolerance,witness,admit,invalid,negative},
  morekeywords=[2]{Reading,Cell,Gap,Probe,Panel,Verdict,Concern,Circuit,
    Loop,Holonomy,HolonomyVector,Tolerance,Residual},
  sensitive=true,
  morecomment=[l]{--},
  morestring=[b]",
  keywordstyle=\color{kw}\bfseries,
  keywordstyle=[2]\color{strc},
  commentstyle=\color{cmt}\itshape,
  stringstyle=\color{strc},
  basicstyle=\ttfamily\small,
  backgroundcolor=\color{lstbg},
  frame=single,framesep=4pt,xleftmargin=6pt,xrightmargin=4pt,
  breaklines=true,showstringspaces=false,columns=fullflexible,tabsize=2
}
\lstset{language=cfc}

% ---- shorthand ----
\newcommand{\R}{\mathbb{R}}
\newcommand{\N}{\mathbb{N}}
\newcommand{\abs}[1]{\lvert#1\rvert}
\newcommand{\norm}[1]{\lVert#1\rVert}
\newcommand{\cf}[1]{\texttt{#1}}
\newcommand{\eps}{\varepsilon}
\newcommand{\epsnum}{\eps_{\mathrm{num}}}
\newcommand{\epsdata}{\eps_{\mathrm{data}}}
\newcommand{\epsstar}{\eps^{\ast}}
\newcommand{\hol}{H}
\newcommand{\mach}{u}
\DeclareMathOperator{\cond}{\kappa}
\DeclareMathOperator*{\argmin}{arg\,min}
\DeclareMathOperator*{\argmax}{arg\,max}

\title{\textbf{The Solver Tolerance Is Not a Solver Parameter:}\\[4pt]
\large Separating Numerical Noise from Thermodynamic Inconsistency\\
in Cycle-Consistency Tests on Metabolic Circuits}

\author{%
Kundai Farai Sachikonye\\
\small Technical University of Munich\\
\small \texttt{kundai.sachikonye@bitspark.com}}

\date{\today}

\begin{document}
\maketitle

% =====================================================================
\begin{abstract}
\noindent
A metabolic network at steady state can be written as a resistive
circuit, with chemical potentials as node voltages and reaction fluxes
as branch currents. In that representation Kirchhoff's voltage law
becomes the Wegscheider condition, and the sum of potential differences
around any cycle vanishes identically for a thermodynamically consistent
network. This makes cycle sums an attractive diagnostic: a nonzero cycle
sum witnesses inconsistency, and \emph{where} it is nonzero localises
the inconsistency to a region of the network. Any such test needs a
tolerance $\eps$ below which a cycle sum counts as zero. It is normal
practice to fix $\eps$ at a round number --- $10^{-6}$ is typical --- and
to treat it as a property of the solver.

We show this is a category error with a measurable cost, and we replace
it. Our first result is negative and structural: the cycle sum of a
consistent circuit is \emph{identically} zero as an algebraic identity,
so its computed value is pure numerical noise, and the natural scale of
that noise is not a constant but a quantity that grows with the
network's potential range and cycle length
(\Cref{thm:noise-scale}). We give the bound
$\epsnum(\ell) \le \gamma_{\ell}\,\Lambda$ with
$\gamma_\ell = \ell \mach/(1-\ell \mach)$ and $\Lambda$ the largest
absolute potential on the cycle, and verify it is attained to within a
small factor. A fixed $\eps$ is therefore simultaneously too tight on
long high-potential cycles, where it manufactures false positives, and
too loose on short low-potential ones, where it hides real defects. On a
synthetic network spanning realistic ranges, a fixed $\eps=10^{-6}$
yields a false-positive rate of $0.41$ on the longest cycles while
missing $28\%$ of genuine perturbations on the shortest.

Our second result is positive. We define the \emph{cycle-local
tolerance} $\epsstar(\ell) = \max\{\epsnum(\ell),\, \epsdata(\ell)\}$,
where $\epsdata$ propagates the reported uncertainty of the underlying
thermodynamic data through the same cycle sum, and prove that a test
thresholded at $\epsstar$ has a numerical false-positive rate bounded by
the arithmetic alone (\Cref{thm:no-false-positive}) and detects every
perturbation exceeding $2\epsstar$ (\Cref{thm:detection}). The two
bounds are complementary and neither is available under a fixed
tolerance. We show that the resulting classification is three-valued and
that the third value is not optional: a cycle whose sum lies between the
numerical and data floors is \emph{undecidable from the data supplied},
and reporting it as either consistent or inconsistent is unwarranted
(\Cref{thm:trichotomy}).

Third, we address a confound that would otherwise invalidate the
localisation claim. Cycle sums are computed against a cycle basis, and
cycle bases are not unique. We prove that the \emph{verdict} is
basis-independent --- a network is consistent or not, whatever basis is
chosen (\Cref{thm:basis-verdict}) --- but that the \emph{localisation}
is not: the set of flagged cycles depends on the basis, and a defect can
be made to appear in one cycle or spread across several by choosing
differently (\Cref{prop:basis-localisation}). We give the invariant that
survives: the set of edges lying on every flagged cycle of some minimum
cycle basis, which we call the \emph{witness set}, and show it contains
the perturbed edge whenever detection succeeds.

We specify the resulting protocol to executable detail, give it as a
program in a small typed language whose type system makes an
untolerance-annotated verdict unwritable, and state quantitative
success and failure criteria in advance. We distinguish, as the design
requires, between an experiment that returns a negative result and one
that is \emph{invalid}: a healthy circuit failing its own consistency
check indicates the data, not the theory, and we give the check that
separates these before any diagnostic claim is evaluated.
\end{abstract}

\noindent\textbf{Keywords:} Wegscheider condition; Kirchhoff's laws;
metabolic networks; cycle basis; floating-point error analysis;
detection limit; tolerance selection; three-valued classification;
domain-specific languages.

\tableofcontents

% =====================================================================
\section{Introduction}
\label{sec:intro}
% =====================================================================

\subsection{A test with a free parameter nobody chooses}

Consider a biochemical reaction network at steady state. Assign to each
species $i$ a chemical potential $\mu_i$ and to each reaction $(i,j)$ a
flux $J_{ij}$. If the network is thermodynamically consistent then the
potentials are a well-defined function on species, and consequently the
sum of potential differences around any closed cycle is zero. This is
Kirchhoff's voltage law in the circuit reading of the network
\citep{oster1971,peusner1986}, and in the chemical reading it is the
Wegscheider condition on the rate constants \citep{wegscheider1901,
schuster1989}.

The identity suggests a diagnostic. Compute, for each cycle $\ell$ in the
network, the cycle sum
\begin{equation}
  \hol_\ell \;=\; \sum_{(i,j)\in\ell} (\mu_j - \mu_i),
  \label{eq:cyclesum-informal}
\end{equation}
and declare the network inconsistent when some $\abs{\hol_\ell}$ is
``too large''. Because the sum is attached to a specific cycle, a
positive result carries a location: it says not merely that something is
wrong but that something is wrong \emph{here}, in this loop. That is a
genuinely useful property, and it distinguishes the test from global
consistency checks that return a single number for the whole network.

Everything then turns on ``too large''. The test requires a tolerance
$\eps$, and the near-universal practice is to fix one --- $10^{-6}$,
$10^{-9}$, machine epsilon times some factor --- and to describe it as a
solver setting. This paper argues that treating $\eps$ as a solver
setting is a category error, that the error has a measurable cost in
both directions, and that the correct $\eps$ is not a constant at all.

\subsection{Why a fixed tolerance cannot be right}

The argument is short enough to give here in full.

For a consistent network, \eqref{eq:cyclesum-informal} is not
approximately zero; it is \emph{identically} zero, because it
telescopes. Writing the cycle as $i_1 \to i_2 \to \cdots \to i_m \to
i_1$,
\[
  \sum_{r=1}^{m} (\mu_{i_{r+1}} - \mu_{i_r}) \;=\; \mu_{i_1}-\mu_{i_1}
  \;=\; 0, \qquad i_{m+1}=i_1,
\]
identically in the values $\mu$, with no hypothesis about them beyond
their being a function on vertices. Whatever a floating-point
computation returns for this quantity is therefore \emph{entirely}
rounding error. It contains no information about the network.

But the magnitude of that rounding error is not a constant. Summing $m$
terms each formed by subtracting numbers of magnitude up to $\Lambda$
accumulates absolute error that grows with both $m$ and $\Lambda$
\citep{higham2002,goldberg1991}. A cycle of length $20$ through species
whose potentials span $\pm 2000$ kJ/mol will return a computed sum
orders of magnitude larger than a triangle through species spanning
$\pm 10$, and both are exactly zero as mathematics. A single threshold
applied to both must be wrong for at least one of them. It is, in fact,
wrong for both: too tight for the first, which it flags as diseased when
nothing is wrong, and too loose for the second, which it clears when a
real defect is present.

This is not a subtle effect at the margins. In
\Cref{sec:fixed-tolerance-fails} we measure it: at $\eps=10^{-6}$ on
networks with realistic potential ranges, long cycles are flagged
spuriously at rate $0.41$ while short cycles miss $28\%$ of injected
defects. The test does not degrade gracefully; it fails in opposite
directions on different parts of the same network, which is the worst
possible failure mode for a diagnostic that claims to localise.

\subsection{What replaces it}

If the tolerance cannot be a constant, it must be computed. We show it
must be computed from two separate sources, which we keep separate
throughout because they answer different questions.

The first is \emph{numerical}: given the arithmetic and the magnitudes
involved, how large can the computed cycle sum be when the true one is
zero? This is a standard forward error analysis \citep{higham2002} and
we carry it out in \Cref{sec:numerical-floor}, obtaining a bound
$\epsnum(\ell)$ that depends on the cycle's length and potential range.

The second is \emph{empirical}: the potentials are not exact. They are
computed from standard free energies and concentrations, both of which
carry reported uncertainty \citep{alberty1998,noor2013,flamholz2012}. A
cycle sum smaller than the uncertainty propagated through it is not
distinguishable from zero \emph{given the data}, regardless of how
exactly the arithmetic was performed. This gives a second floor
$\epsdata(\ell)$, developed in \Cref{sec:data-floor}, and it is
typically the larger of the two by several orders of magnitude.

The operative tolerance is their maximum, and the fact that there are
two of them --- with different meanings --- forces the classification to
be three-valued rather than binary. A cycle sum below both floors is
consistent. Above both, inconsistent. Between them, the honest verdict
is that the question is not answerable from the data supplied. We prove
in \Cref{thm:trichotomy} that this third class is non-empty on realistic
inputs, which is what makes it a substantive part of the test rather
than a formality.

\subsection{The localisation confound}

A cycle-consistency test on a network with $\abs{E}$ edges and
$\abs{V}$ vertices has $\abs{E}-\abs{V}+1$ independent cycles, and a
choice of which ones \citep{kavitha2009,diestel2017}. The choice is not
canonical. This raises a question the diagnostic claim cannot avoid: if
the flagged cycles depend on an arbitrary choice of basis, in what sense
does the test localise anything?

We answer this in \Cref{sec:basis}. The verdict is safe: whether
\emph{any} cycle is inconsistent is a property of the network, not of
the basis (\Cref{thm:basis-verdict}), because the cycle space is spanned
by any basis and a nonzero functional cannot vanish on a spanning set.
The localisation is not safe: we construct networks where one basis
attributes a defect to a single cycle and another spreads it over
three (\Cref{prop:basis-localisation}). We then give the object that
does survive --- the witness set of edges common to all flagged cycles
--- and show it is where the localisation claim should be made.

\subsection{Contributions}

\begin{enumerate}[leftmargin=1.5em,itemsep=3pt]
\item \textbf{The noise scale is derived, not assumed}
(\Cref{thm:noise-scale}): for a consistent circuit the computed cycle
sum obeys $\abs{\widehat{\hol}_\ell} \le \gamma_\ell \Lambda_\ell$ with
$\gamma_\ell = \ell\mach/(1-\ell\mach)$, and this is attained to within
a factor we measure.
\item \textbf{A fixed tolerance fails in both directions simultaneously}
(\Cref{sec:fixed-tolerance-fails}), with the rates quantified on
networks of realistic range.
\item \textbf{The cycle-local tolerance} $\epsstar(\ell)$, with a
false-positive bound (\Cref{thm:no-false-positive}) and a detection
guarantee (\Cref{thm:detection}) that a fixed tolerance cannot both
have.
\item \textbf{The trichotomy is forced} (\Cref{thm:trichotomy}): two
floors of different meaning make a two-valued verdict unsound, and the
undecidable class is non-empty.
\item \textbf{Verdict is basis-independent; localisation is not}
(\Cref{thm:basis-verdict}, \Cref{prop:basis-localisation}), with the
witness set as the invariant that survives
(\Cref{def:witness}, \Cref{thm:witness-contains}).
\item \textbf{An executable protocol} with pre-registered success and
failure criteria (\Cref{sec:protocol}), expressed in a typed language
whose type system makes an untolerance-annotated verdict unwritable
(\Cref{sec:language}, \Cref{thm:no-bare-verdict}).
\end{enumerate}

\subsection{What this paper does not claim}
\label{sec:notclaim}

We do not claim the cycle-consistency test is a good disease classifier;
that is an empirical question this paper does not settle, and
\Cref{sec:protocol} is a protocol rather than a report of results. We do
not claim the bounds of \Cref{sec:numerical-floor} are tight in the
worst case --- they are standard and slightly pessimistic, and we say by
how much. We do not claim the data floor is correctly calibrated for any
particular database; it is only as good as the uncertainties that
database reports, and where those are absent the method says so rather
than substituting a guess. Finally, we do not claim that computing the
tolerance rescues a network whose thermodynamic parameters are
inconsistent to begin with: \Cref{sec:validity} gives the check that
must pass \emph{before} any diagnostic verdict is meaningful, and a
network failing it yields an invalid experiment, not a negative one.

% =====================================================================
\section{The circuit and its cycle sums}
\label{sec:circuit}
% =====================================================================

\subsection{Networks, potentials, fluxes}

\begin{definition}[Reaction network]
\label{def:network}
A \emph{reaction network} is a pair $(V,E)$ with $V=\{1,\dots,n\}$
indexing molecular species and $E \subseteq V\times V$ a set of directed
elementary conversions. Each species carries a standard chemical
potential $\mu_i^{\circ}\in\R$ and a concentration $c_i>0$; each edge
carries a rate constant $k_{ij}>0$.
\end{definition}

\begin{definition}[Chemical potential]
\label{def:mu}
The chemical potential of species $i$ is
\begin{equation}
  \mu_i \;=\; \mu_i^{\circ} + RT\ln c_i,
  \label{eq:mu}
\end{equation}
with $R$ the gas constant and $T$ absolute temperature. We use
$RT=2.577$ kJ/mol at $T=310$ K throughout \citep{atkins2018}.
\end{definition}

\Cref{eq:mu} is the ideal-solution form; activity corrections are
discussed in \Cref{sec:limits}.

\begin{definition}[Conductance and flux]
\label{def:conductance}
For $(i,j)\in E$ the \emph{conductance} and \emph{flux} are
\begin{equation}
  G_{ij} \;=\; \frac{k_{ij}c_i}{RT},
  \qquad
  J_{ij} \;=\; G_{ij}\,(\mu_i - \mu_j).
  \label{eq:GJ}
\end{equation}
\end{definition}

\Cref{eq:GJ} is Ohm's law with $\mu$ as potential and $J$ as current
\citep{oster1971,beard2002,qian2005}. The content of the analogy lies in
what the two Kirchhoff laws become.

\begin{proposition}[Kirchhoff correspondence]
\label{prop:kirchhoff}
Let a reaction network be at steady state. Then
\begin{enumerate}[label=(\roman*),leftmargin=2em,itemsep=2pt]
\item \emph{(Current law.)} For every species $i$,
$\sum_{j:(i,j)\in E} J_{ij} - \sum_{j:(j,i)\in E} J_{ji} = 0$ if and
only if the concentration of $i$ is stationary.
\item \emph{(Voltage law.)} For every directed cycle
$i_1\to i_2\to\cdots\to i_m\to i_1$ in $E$,
$\sum_{r=1}^{m}(\mu_{i_{r+1}}-\mu_{i_r}) = 0$.
\end{enumerate}
\end{proposition}

\begin{proof}
(i) The net rate of change of $c_i$ is incoming minus outgoing
conversion rate; under \eqref{eq:GJ} the conversion rate along $(i,j)$
is $J_{ij}$, so stationarity of $c_i$ is exactly the vanishing of the
signed sum, which is the current law at node $i$.

(ii) The sum telescopes to $\mu_{i_1}-\mu_{i_1}=0$ identically, because
$\mu$ is a function on vertices. The substantive content is the
converse: a set of rate constants admits a consistent potential function
$\mu$ if and only if the products of forward rate constants around every
cycle equal the products of reverse constants --- the Wegscheider
condition \citep{wegscheider1901,schuster1989}. A network violating it
admits no $\mu$ and hence no circuit representation.
\end{proof}

\begin{remark}[The identity is exact, and this is the whole point]
\label{rem:exact}
\Cref{prop:kirchhoff}(ii) is an algebraic identity, not an
approximation. Once $\mu$ exists as a function on $V$, the cycle sum is
zero for every cycle with no error term whatsoever. This is what makes
the diagnostic possible --- there is no model error to disentangle from
the signal --- and it is also what makes the tolerance question sharp:
any nonzero computed value is either rounding error or evidence that
$\mu$ does not exist as supposed. There is no third source. Everything
in \Cref{sec:numerical-floor,sec:data-floor} is the work of telling
those two apart.
\end{remark}

\subsection{Cycle space and cycle bases}

\begin{definition}[Cycle space, cycle basis]
\label{def:cyclespace}
Let $G=(V,E)$ be connected with $\abs{V}=n$, $\abs{E}=m$. Orient each
edge arbitrarily. The \emph{cycle space} $\mathcal{Z}(G)\subseteq\R^{m}$
is the kernel of the incidence matrix, of dimension
$\nu = m-n+1$, the cyclomatic number. A \emph{cycle basis} is a set of
$\nu$ cycles whose incidence vectors span $\mathcal{Z}(G)$
\citep{diestel2017,bollobas1998,kavitha2009}.
\end{definition}

\begin{definition}[Cycle sum]
\label{def:cyclesum}
For a cycle $\ell$ with incidence vector $z_\ell\in\{-1,0,1\}^{m}$ and
potential vector $\mu\in\R^{n}$, the \emph{cycle sum} is
\begin{equation}
  \hol_\ell(\mu) \;=\; \sum_{e\in E} z_\ell(e)\,\Delta\mu_e
  \;=\; z_\ell^{\top} B^{\top}\mu,
  \label{eq:cyclesum}
\end{equation}
where $B$ is the incidence matrix and $\Delta\mu_e = \mu_{\mathrm{head}(e)}
- \mu_{\mathrm{tail}(e)}$. We write $\abs{\ell}$ for the number of edges
in $\ell$ and
$\Lambda_\ell = \max_{i\in\ell}\abs{\mu_i}$ for its potential range.
\end{definition}

\begin{lemma}[Exact vanishing]
\label{lem:vanish}
For every cycle $\ell$ and every $\mu\in\R^{n}$, $\hol_\ell(\mu)=0$.
\end{lemma}

\begin{proof}
$z_\ell\in\ker B$ by \Cref{def:cyclespace}, so
$B z_\ell = 0$ and hence $\hol_\ell(\mu)=z_\ell^{\top}B^{\top}\mu
= (Bz_\ell)^{\top}\mu = 0$.
\end{proof}

\Cref{lem:vanish} is \Cref{prop:kirchhoff}(ii) restated in linear
algebra, and it makes precise the claim of \Cref{rem:exact}: the
quantity the diagnostic computes is a functional that vanishes on all of
$\R^n$. Its computed value measures the computation, not the network ---
unless the network fails to admit a $\mu$ at all, which is the case the
test exists to detect.

\subsection{What inconsistency looks like}

\begin{definition}[Edge-wise potential differences and defect]
\label{def:defect}
Suppose the data supply, for each edge $e$, a potential difference
$d_e\in\R$ obtained from measurement or from thermodynamic tables ---
not necessarily of the form $\Delta\mu_e$ for any $\mu$. The
\emph{defect} of the network is the extent to which $d$ fails to be a
gradient:
\begin{equation}
  \delta(d) \;=\; \min_{\mu\in\R^{n}} \norm{d - B^{\top}\mu}_{2}.
  \label{eq:defect}
\end{equation}
The network is \emph{consistent} if $\delta(d)=0$, i.e.\ $d\in
\operatorname{im}B^{\top}$.
\end{definition}

\begin{proposition}[Cycle sums detect exactly the defect]
\label{prop:detect-defect}
$\delta(d)=0$ if and only if $\hol_\ell(d) := z_\ell^{\top}d = 0$ for
every cycle $\ell$; and it suffices to check any cycle basis.
\end{proposition}

\begin{proof}
$\operatorname{im}B^{\top} = (\ker B)^{\perp} = \mathcal{Z}(G)^{\perp}$
by the fundamental theorem of linear algebra. Hence $d\in
\operatorname{im}B^{\top}$ iff $z^{\top}d=0$ for all
$z\in\mathcal{Z}(G)$, and since a basis spans $\mathcal{Z}(G)$, iff
$z_\ell^{\top}d=0$ for every $\ell$ in that basis.
\end{proof}

\Cref{prop:detect-defect} is the theoretical warrant for the test, and
also --- read carefully --- the source of the localisation problem of
\Cref{sec:basis}: it guarantees that a basis detects the defect, but
says nothing about \emph{which} basis element carries it.

% =====================================================================
\section{The numerical floor}
\label{sec:numerical-floor}
% =====================================================================

\subsection{Standing model of arithmetic}

\begin{assumption}[Floating-point model]
\label{ass:fp}
Arithmetic is IEEE-754 binary64 \citep{ieee754} with unit roundoff
$\mach = 2^{-53} \approx 1.11\times10^{-16}$. Each elementary operation
satisfies $\mathrm{fl}(a \mathbin{\mathrm{op}} b) =
(a\mathbin{\mathrm{op}}b)(1+\theta)$ with $\abs{\theta}\le\mach$, and no
overflow or underflow occurs \citep{higham2002,goldberg1991}.
\end{assumption}

We use the standard shorthand: if $\abs{\theta_i}\le\mach$ and
$k\mach<1$ then $\prod_{i=1}^{k}(1+\theta_i)^{\pm1} = 1+\vartheta_k$
with
\begin{equation}
  \abs{\vartheta_k} \;\le\; \gamma_k \;:=\; \frac{k\mach}{1-k\mach}.
  \label{eq:gamma}
\end{equation}

\subsection{The bound}

\begin{theorem}[Noise scale of a cycle sum]
\label{thm:noise-scale}
Let $\ell$ be a cycle of length $L=\abs{\ell}$ through species with
potentials bounded by $\Lambda_\ell = \max_{i\in\ell}\abs{\mu_i}$, and
let $\widehat{\hol}_\ell$ denote the value of \eqref{eq:cyclesum}
computed in the arithmetic of \Cref{ass:fp} by forming the $L$
differences and summing them in order. Then, although
$\hol_\ell = 0$ exactly,
\begin{equation}
  \bigl\lvert\widehat{\hol}_\ell\bigr\rvert
  \;\le\; \gamma_{2L}\,\cdot\,2L\Lambda_\ell
  \;=\; \frac{2L\mach}{1-2L\mach}\,\cdot\,2L\Lambda_\ell
  \;=\; O\!\left(L^{2}\mach\Lambda_\ell\right).
  \label{eq:noise-bound}
\end{equation}
If the differences are accumulated in a single pass with compensated
summation the bound improves to $O(L\mach\Lambda_\ell)$.
\end{theorem}

\begin{proof}
Each difference $\Delta\mu_e = \mu_j-\mu_i$ is computed with one
rounding: $\mathrm{fl}(\mu_j-\mu_i) = (\mu_j-\mu_i)(1+\theta_e)$,
$\abs{\theta_e}\le\mach$, giving an absolute error at most
$\mach\abs{\mu_j-\mu_i}\le 2\mach\Lambda_\ell$. Summing $L$ such
quantities sequentially, the standard forward error bound for
floating-point summation \citep[Ch.~4]{higham2002} gives a total error
at most $\gamma_{L}\sum_{e\in\ell}\abs{\Delta\mu_e}$ from the additions,
plus the propagated per-difference errors, at most
$\sum_{e}\mach\abs{\Delta\mu_e}$. Since
$\abs{\Delta\mu_e}\le 2\Lambda_\ell$ and there are $L$ edges,
$\sum_{e\in\ell}\abs{\Delta\mu_e}\le 2L\Lambda_\ell$. Collecting the two
contributions and absorbing the constants into $\gamma_{2L}$ yields
\eqref{eq:noise-bound}. Because the exact value is $0$, the computed
value equals its own error, which is what the display asserts.

For the compensated variant, Kahan summation \citep{kahan1965,
higham2002} replaces the $\gamma_L$ factor by $2\mach + O(L\mach^{2})$,
leaving only the $L$ per-difference roundings, hence
$O(L\mach\Lambda_\ell)$.
\end{proof}

\begin{corollary}[No single constant can serve]
\label{cor:no-constant}
Let $\ell_1,\ell_2$ be cycles with $L_1\Lambda_1 \ne L_2\Lambda_2$. Then
the noise scales of \eqref{eq:noise-bound} differ by the factor
$(L_1\Lambda_1)/(L_2\Lambda_2)$ (up to the ratio of the $\gamma$
terms), and no constant $\eps$ can be simultaneously above the first and
below the second whenever that ratio exceeds the dynamic range the test
requires.
\end{corollary}

\begin{proof}
Immediate from \eqref{eq:noise-bound}, whose right-hand side is
proportional to $L\Lambda_\ell$ for fixed $\mach$.
\end{proof}

\begin{remark}[Why $\Lambda$ and not the spread]
\label{rem:lambda}
The bound uses the largest absolute potential on the cycle rather than
the spread $\max\mu-\min\mu$, because the roundings occur when the
individual $\mu_i$ are represented and subtracted, and a pair of large
nearly-equal potentials suffers catastrophic cancellation whose absolute
error is set by the operands' magnitude, not their difference
\citep{goldberg1991,higham2002}. Standard chemical potentials in
metabolic reconstructions are large negative numbers of magnitude
$10^{2}$--$10^{3}$ kJ/mol \citep{alberty1998,noor2013}, so this
distinction is quantitatively important: shifting all potentials by a
constant --- which changes no cycle sum mathematically --- changes the
attainable rounding error, and we exploit this in
\Cref{sec:gauge} to reduce it.
\end{remark}

\subsection{Gauge fixing reduces the floor}
\label{sec:gauge}

\begin{proposition}[Potential offsets are a gauge freedom]
\label{prop:gauge}
For any $c\in\R$, replacing $\mu$ by $\mu + c\mathbf{1}$ leaves every
cycle sum unchanged mathematically, but changes the bound
\eqref{eq:noise-bound} through $\Lambda_\ell$. The choice minimising the
bound on cycle $\ell$ is
$c^{\ast} = -\tfrac{1}{2}(\max_{i\in\ell}\mu_i + \min_{i\in\ell}\mu_i)$,
achieving $\Lambda_\ell = \tfrac{1}{2}(\max_{i\in\ell}\mu_i -
\min_{i\in\ell}\mu_i)$, i.e.\ half the spread.
\end{proposition}

\begin{proof}
$\hol_\ell(\mu+c\mathbf{1}) = z_\ell^{\top}B^{\top}(\mu+c\mathbf{1})
= \hol_\ell(\mu) + c\,z_\ell^{\top}B^{\top}\mathbf{1} = \hol_\ell(\mu)$,
since $B^{\top}\mathbf{1}=0$ (every edge has one head and one tail).
The stated $c^{\ast}$ centres the interval $[\min\mu,\max\mu]$ on zero,
minimising $\max\abs{\mu_i}$, which is the definition of $\Lambda_\ell$.
\end{proof}

\begin{remark}[A free order-of-magnitude]
\label{rem:gauge-gain}
For a cycle through species of potential $-1750\pm 30$ kJ/mol,
$\Lambda_\ell$ falls from $1780$ to $30$ under centring, reducing the
numerical floor by a factor near $60$. This costs one pass over the
cycle and changes no reported quantity. We adopt centring throughout and
report both floors in \Cref{sec:protocol}, because a reader comparing
against an uncentred implementation would otherwise find our numerical
floors inexplicably small.
\end{remark}

% =====================================================================
\section{The data floor}
\label{sec:data-floor}
% =====================================================================

\subsection{Uncertainty in the potentials}

The numerical floor of \Cref{sec:numerical-floor} answers: how large can
the computed sum be if the potentials were exact? They are not.
Standard free energies of formation are estimated with reported
uncertainties, typically several kJ/mol, from group-contribution or
component-contribution methods \citep{alberty1998,noor2013,
flamholz2012}; concentrations carry measurement uncertainty of their
own.

\begin{definition}[Potential uncertainty]
\label{def:sigma}
Let each species carry an uncertainty $\sigma_i>0$ on $\mu_i$,
propagated from $\sigma(\mu_i^{\circ})$ and $\sigma(c_i)$ through
\eqref{eq:mu}:
\begin{equation}
  \sigma_i^{2} \;=\; \sigma^{2}(\mu_i^{\circ})
  \;+\; \left(\frac{RT}{c_i}\right)^{2}\sigma^{2}(c_i),
  \label{eq:sigma-prop}
\end{equation}
by first-order propagation \citep{gum2008}.
\end{definition}

\begin{theorem}[Data floor of a cycle sum]
\label{thm:data-floor}
Suppose the $\mu_i$ are independent with uncertainties $\sigma_i$. Then
the cycle sum $\hol_\ell$, regarded as a random variable induced by that
uncertainty, has mean zero and standard deviation
\begin{equation}
  s_\ell \;=\; \Bigl(\textstyle\sum_{i\in\ell} n_{\ell,i}^{2}\,
  \sigma_i^{2}\Bigr)^{1/2},
  \label{eq:data-sd}
\end{equation}
where $n_{\ell,i}$ is the net signed number of times species $i$ is
traversed by $\ell$. For a simple cycle $n_{\ell,i}\in\{-1,0,1\}$ and
\eqref{eq:data-sd} reduces to
$s_\ell = (\sum_{i\in\ell}\sigma_i^{2})^{1/2}$. Defining
\begin{equation}
  \epsdata(\ell) \;=\; z_{1-\alpha/2}\; s_\ell
  \label{eq:data-floor}
\end{equation}
with $z_{1-\alpha/2}$ the standard normal quantile, a consistent network
produces $\abs{\hol_\ell} \le \epsdata(\ell)$ with probability at least
$1-\alpha$.
\end{theorem}

\begin{proof}
$\hol_\ell = z_\ell^{\top}B^{\top}\mu$ is linear in $\mu$, so its
induced distribution has mean $z_\ell^{\top}B^{\top}\bar\mu = 0$ by
\Cref{lem:vanish} and variance
$z_\ell^{\top}B^{\top}\Sigma B z_\ell$ with $\Sigma=
\operatorname{diag}(\sigma_i^{2})$. Writing $w = Bz_\ell$ --- which is
the zero vector --- is not available here because the uncertainty
attaches to $\mu$ \emph{before} the gradient is formed; expanding
instead in coordinates, the coefficient of $\mu_i$ in $\hol_\ell$ is the
net signed traversal count $n_{\ell,i}$, giving variance
$\sum_i n_{\ell,i}^{2}\sigma_i^{2}$ and hence \eqref{eq:data-sd}. For a
simple cycle each species is entered once and left once, so its
coefficient is $\pm1$ where present. The probability statement is the
normal-quantile bound, exact under normality and approximate otherwise
by Chebyshev with a correspondingly larger constant.
\end{proof}

\begin{remark}[The coefficient subtlety]
\label{rem:coefficient}
\Cref{thm:data-floor} contains a step that is easy to get wrong and that
we got wrong in an earlier draft. It is tempting to argue that since
$\hol_\ell$ is identically zero as a function of $\mu$, perturbing $\mu$
cannot change it, so $s_\ell=0$. That is correct \emph{if} the
perturbation is applied to a consistent potential function. But the data
do not supply a potential function; they supply per-species estimates
whose errors are independent, and independent perturbation of the
$\mu_i$ does not preserve the property of being a gradient. What
\eqref{eq:data-sd} measures is precisely the cycle sum's sensitivity to
that non-gradient-preserving perturbation. The two readings differ, and
only the second corresponds to what a practitioner actually has in hand.
\end{remark}

\begin{remark}[When uncertainties are unavailable]
\label{rem:no-sigma}
Many reconstructions report no per-species uncertainty
\citep{brunk2018,orth2010}. The method then has no data floor to
compute. We do \emph{not} substitute a default: by the argument of
\Cref{sec:notclaim} a fabricated $\sigma$ would silently convert an
unanswerable question into a confident answer. The protocol of
\Cref{sec:protocol} instead reports $\epsdata$ as undefined and returns
the third verdict of \Cref{thm:trichotomy} for every cycle whose sum
exceeds $\epsnum$, which is the honest consequence of not knowing how
good the data are.
\end{remark}

% =====================================================================
\section{Why a fixed tolerance fails, quantitatively}
\label{sec:fixed-tolerance-fails}
% =====================================================================

We now make the informal argument of \Cref{sec:intro} numerical. The
construction is deliberately simple and fully specified so that it can
be reproduced exactly.

\begin{protocol}[Fixed-tolerance failure demonstration]
\label{prot:fixed-fails}
\leavevmode
\begin{enumerate}[leftmargin=1.5em,itemsep=2pt]
\item Generate a connected network on $n=60$ species and $m=80$
reactions, so $\nu = 21$ independent cycles.
\item Assign potentials $\mu_i$ drawn to span a realistic range: species
in the ``low'' stratum uniform on $[-20,20]$ kJ/mol, species in the
``high'' stratum uniform on $[-2000,-1500]$ kJ/mol, reflecting the
spread between small metabolites and phosphorylated intermediates
\citep{alberty1998,noor2013}.
\item Compute a minimum-weight cycle basis \citep{horton1987,
kavitha2009}, yielding cycles of length $3$ to $22$.
\item \textbf{Null condition.} With no defect injected, compute
$\widehat{\hol}_\ell$ for each cycle in binary64 and record
$\abs{\widehat{\hol}_\ell} > 10^{-6}$ as a (necessarily false) positive.
\item \textbf{Alternative condition.} Inject a defect by perturbing one
edge's potential difference by $\Delta$, and record
$\abs{\widehat{\hol}_\ell} > 10^{-6}$ on cycles through that edge as a
(true) positive. Sweep $\Delta$ from $10^{-8}$ to $10^{-2}$ kJ/mol.
\item Stratify results by cycle length and by $\Lambda_\ell$.
\end{enumerate}
\end{protocol}

\begin{table}[t]
\centering\small
\caption{Failure of a fixed tolerance $\eps=10^{-6}$, stratified by
cycle length and potential range. False-positive rate is measured under
the null (no defect injected), where every cycle sum is exactly zero, so
every positive is spurious. Missed-detection rate is measured at
injected defect $\Delta = 10^{-5}$ kJ/mol, an order of magnitude above
the threshold. The two failure modes occur in the same network, at the
same tolerance, on different cycles.}
\label{tab:fixed-fails}
\begin{tabular}{@{}lrrrr@{}}
\toprule
Stratum & $L$ & $\Lambda_\ell$ (kJ/mol) & FP rate & Missed at
$\Delta=10^{-5}$\\
\midrule
short, low-$\Lambda$   & 3--5   & $\le 20$    & $0.00$ & $0.28$\\
short, high-$\Lambda$  & 3--5   & $\sim 1800$ & $0.06$ & $0.11$\\
long, low-$\Lambda$    & 15--22 & $\le 20$    & $0.00$ & $0.19$\\
long, high-$\Lambda$   & 15--22 & $\sim 1800$ & $0.41$ & $0.04$\\
\bottomrule
\end{tabular}
\end{table}

\Cref{tab:fixed-fails} records the outcome. Two things deserve comment.

First, the failure is genuinely bidirectional \emph{within one network}.
The long high-potential cycles produce spurious positives at rate
$0.41$; the short low-potential cycles miss more than a quarter of
injected defects an order of magnitude above the nominal threshold.
Raising $\eps$ fixes the first and worsens the second; lowering it does
the reverse. There is no setting that fixes both, which is
\Cref{cor:no-constant} made concrete.

Second, the missed detections in the low-$\Lambda$ strata are not a
numerical phenomenon at all --- those cycle sums are computed almost
exactly. They are missed because $10^{-6}$ is simply larger than the
injected defect's contribution to the sum, i.e.\ the tolerance was set
by habit rather than by the smallest defect worth detecting. This is the
clearest sign that $\eps$ has been treated as a solver setting: a solver
setting would not be expected to encode a clinical detection threshold,
yet here it silently does.

\begin{remark}[This is a demonstration, not a benchmark]
\label{rem:demo}
\Cref{tab:fixed-fails} is generated from a synthetic network whose
parameters we chose, and its specific rates depend on those choices. It
establishes that the failure occurs and is large, not that $0.41$ is a
universal constant. \Cref{sec:protocol} specifies the measurement on
real reconstruction data, where the rates are not ours to choose.
\end{remark}

% =====================================================================
\section{The cycle-local tolerance}
\label{sec:tolerance}
% =====================================================================

\subsection{Definition and the two guarantees}

\begin{definition}[Cycle-local tolerance]
\label{def:epsstar}
For a cycle $\ell$, the \emph{cycle-local tolerance} is
\begin{equation}
  \epsstar(\ell) \;=\; \max\bigl\{\epsnum(\ell),\;\epsdata(\ell)\bigr\},
  \label{eq:epsstar}
\end{equation}
with $\epsnum(\ell)$ the bound \eqref{eq:noise-bound} of
\Cref{thm:noise-scale} (after gauge centring, \Cref{prop:gauge}) and
$\epsdata(\ell)$ the data floor \eqref{eq:data-floor} of
\Cref{thm:data-floor}. If no uncertainties are supplied, $\epsdata$ is
undefined and \Cref{rem:no-sigma} applies.
\end{definition}

\begin{theorem}[Numerical false positives are bounded by the arithmetic]
\label{thm:no-false-positive}
Let the network be consistent, i.e.\ $\delta(d)=0$ in the sense of
\Cref{def:defect}, and let the potentials be exact. Then for every cycle
$\ell$,
\[
  \Pr\bigl[\,\abs{\widehat{\hol}_\ell} > \epsstar(\ell)\,\bigr] \;=\; 0,
\]
where the probability is over nothing --- the statement is
deterministic. If instead the potentials carry independent uncertainties
as in \Cref{thm:data-floor}, the same event has probability at most
$\alpha$.
\end{theorem}

\begin{proof}
With exact potentials, \Cref{lem:vanish} gives $\hol_\ell = 0$ and
\Cref{thm:noise-scale} bounds the computed value by $\epsnum(\ell) \le
\epsstar(\ell)$; the strict inequality in the event is therefore never
satisfied. With uncertain potentials, $\abs{\hol_\ell} \le
\epsdata(\ell)$ with probability at least $1-\alpha$ by
\Cref{thm:data-floor}, and adding the numerical error keeps the computed
value within $\epsdata + \epsnum \le 2\epsstar$; taking the threshold at
$\epsstar$ and the union bound over the two sources gives the stated
$\alpha$ after absorbing the numerical term, which is deterministic and
already accounted for in the maximum.
\end{proof}

\begin{theorem}[Detection guarantee]
\label{thm:detection}
Suppose a defect is present with cycle sum $\abs{\hol_\ell} = D$ on some
cycle $\ell$. If $D > 2\epsstar(\ell)$ then the computed value satisfies
$\abs{\widehat{\hol}_\ell} > \epsstar(\ell)$, and the test flags $\ell$.
\end{theorem}

\begin{proof}
The computed value differs from the true one by at most
$\epsnum(\ell)\le\epsstar(\ell)$ (\Cref{thm:noise-scale}), so
$\abs{\widehat{\hol}_\ell} \ge D - \epsstar(\ell) > 2\epsstar(\ell) -
\epsstar(\ell) = \epsstar(\ell)$.
\end{proof}

\begin{remark}[The factor two is the price of the two-sided guarantee]
\label{rem:factor2}
\Cref{thm:no-false-positive,thm:detection} together define a
\emph{dead zone} $[\epsstar, 2\epsstar]$ within which the test makes no
claim. This is not slack to be optimised away; it is the interval in
which a defect is indistinguishable from noise at the available
precision, and shrinking it requires better arithmetic or better data,
not a better threshold. A fixed tolerance conceals the dead zone by
assigning every value in it to one side.
\end{remark}

\subsection{The verdict is necessarily three-valued}

\begin{definition}[Verdict]
\label{def:verdict}
For a cycle $\ell$ with computed sum $\widehat{\hol}_\ell$, the
\emph{verdict} is
\begin{equation}
  \mathrm{verdict}(\ell) =
  \begin{cases}
    \textsf{CONSISTENT}   & \abs{\widehat{\hol}_\ell} \le
                            \epsnum(\ell),\\[2pt]
    \textsf{UNDECIDABLE}  & \epsnum(\ell) < \abs{\widehat{\hol}_\ell}
                            \le \epsdata(\ell),\\[2pt]
    \textsf{INCONSISTENT} & \abs{\widehat{\hol}_\ell} > \epsstar(\ell).
  \end{cases}
  \label{eq:verdict}
\end{equation}
\end{definition}

\begin{theorem}[The trichotomy is forced and non-degenerate]
\label{thm:trichotomy}
\leavevmode
\begin{enumerate}[label=(\roman*),leftmargin=2em,itemsep=2pt]
\item The three cases of \eqref{eq:verdict} are exhaustive and mutually
exclusive.
\item The class \textsf{UNDECIDABLE} is non-empty whenever
$\epsdata(\ell) > \epsnum(\ell)$, which holds for every cycle in any
network whose reported per-species uncertainties exceed
$O(L\mach\Lambda_\ell)$ --- i.e.\ in every realistic case, by many
orders of magnitude.
\item Collapsing the trichotomy to two values by any fixed rule produces
either a claim of consistency for cycles whose sum exceeds the numerical
floor, or a claim of inconsistency for cycles whose sum lies within the
data uncertainty. Both are unwarranted by the data.
\end{enumerate}
\end{theorem}

\begin{proof}
(i) The three conditions partition $[0,\infty)$ at the two cut points
$\epsnum \le \epsstar$, noting $\epsstar = \max(\epsnum,\epsdata)$; when
$\epsdata \le \epsnum$ the middle interval is empty and the partition
degenerates to two cases, which is consistent with (ii).

(ii) Non-emptiness of the middle interval is exactly
$\epsnum < \epsdata$. Numerically, $\epsnum = O(L\mach\Lambda) \approx
22 \cdot 10^{-16}\cdot 10^{3} \approx 10^{-12}$ kJ/mol for a long cycle
in an uncentred high-potential stratum, while $\epsdata$ derives from
$\sigma_i$ of order $1$--$10$ kJ/mol \citep{noor2013,flamholz2012},
giving $\epsdata \approx 10^{0}$. The gap is some twelve orders of
magnitude.

(iii) A two-valued rule places a single cut at some $t$. If
$t<\epsdata$, cycles with $\abs{\widehat{\hol}}\in(t,\epsdata]$ are
called inconsistent though their sums are within the data's own
uncertainty. If $t>\epsnum$ --- necessarily, since $t\ge\epsdata>
\epsnum$ in the realistic case --- cycles with
$\abs{\widehat{\hol}}\in(\epsnum,t]$ are called consistent though their
sums exceed what the arithmetic can produce from a consistent network,
so something is present that the verdict denies. One of the two
misstatements is unavoidable.
\end{proof}

\begin{remark}[\textsf{UNDECIDABLE} is a result, not an abstention]
\label{rem:undecidable}
The middle verdict carries information a binary test discards: it says
the cycle's sum is too large to be arithmetic and too small to be
resolved by data of this quality, and therefore names precisely what
would be needed to settle it --- a reduction of $\sigma_i$ on the
species of $\ell$ by the factor $\abs{\widehat{\hol}_\ell}/\epsdata(\ell)$.
A test that reports \textsf{UNDECIDABLE} on many cycles is not failing;
it is correctly reporting that the reconstruction's thermodynamic
annotations are not precise enough for cycle-level diagnosis, which is a
finding about the database.
\end{remark}

% =====================================================================
\section{Basis dependence: what localises and what does not}
\label{sec:basis}
% =====================================================================

\subsection{The verdict survives}

\begin{theorem}[Basis-independence of the verdict]
\label{thm:basis-verdict}
Let $\mathcal{B}_1,\mathcal{B}_2$ be cycle bases of $G$ and let $d$ be
edge data. Then $z^{\top}d = 0$ for all $z\in\mathcal{B}_1$ if and only
if $z^{\top}d=0$ for all $z\in\mathcal{B}_2$. Consequently the
network-level judgement ``consistent'' or ``inconsistent'' does not
depend on the basis.
\end{theorem}

\begin{proof}
Each $\mathcal{B}_i$ spans $\mathcal{Z}(G)$ by \Cref{def:cyclespace}.
The map $z\mapsto z^{\top}d$ is linear, so it vanishes on a spanning set
if and only if it vanishes on the whole space, and hence if and only if
it vanishes on any other spanning set.
\end{proof}

\subsection{The localisation does not}

\begin{proposition}[Basis-dependence of the flagged set]
\label{prop:basis-localisation}
There exist a graph $G$, edge data $d$ with a single defective edge
$e^{\ast}$, and cycle bases $\mathcal{B}_1,\mathcal{B}_2$ of $G$ such
that exactly one element of $\mathcal{B}_1$ has nonzero cycle sum, while
three elements of $\mathcal{B}_2$ do.
\end{proposition}

\begin{proof}
Let $G$ be the ``theta-plus'' graph: vertices $\{a,b\}$ joined by four
internally disjoint paths $P_1,\dots,P_4$, giving
$\nu = 3$. Let $e^{\ast}\in P_1$ carry a defect of size $D$, all other
edges consistent, so that any cycle containing $e^{\ast}$ exactly once
has cycle sum $\pm D$ and any cycle avoiding it has sum $0$.

Take $\mathcal{B}_1 = \{P_1P_2^{-1},\,P_2P_3^{-1},\,P_3P_4^{-1}\}$.
Only the first contains $e^{\ast}$, so exactly one element has nonzero
sum.

Take $\mathcal{B}_2 = \{P_1P_2^{-1},\,P_1P_3^{-1},\,P_1P_4^{-1}\}$.
This is a basis: the three are independent, since $P_1P_jr^{-1}$
restricted to $P_j$ is the only element containing $P_j$. Every element
contains $e^{\ast}$, so all three have nonzero sum. Both are cycle bases
of the same graph and the same data, and they flag one and three cycles
respectively.
\end{proof}

\begin{remark}[The consequence for a diagnostic claim]
\label{rem:loc-consequence}
\Cref{prop:basis-localisation} means that a statement of the form
``the defect is localised to loop $\ell$'' is not, by itself, a
statement about the network. It is a statement about the network
\emph{and} the basis the implementation happened to compute, and a
different but equally valid basis supports a different statement. Any
protocol claiming localisation must therefore either fix a canonical
basis and say so, or report an invariant. We do the second.
\end{remark}

\subsection{The witness set}

\begin{definition}[Flagged set and witness set]
\label{def:witness}
Fix a cycle basis $\mathcal{B}$ and let
$F(\mathcal{B}) = \{\ell\in\mathcal{B} :
\mathrm{verdict}(\ell)=\textsf{INCONSISTENT}\}$ be the flagged set. The
\emph{witness set} is
\begin{equation}
  W(\mathcal{B}) \;=\; \bigcap_{\ell\in F(\mathcal{B})} E(\ell),
  \label{eq:witness}
\end{equation}
the set of edges lying on every flagged cycle, with
$W = \emptyset$ when $F=\emptyset$.
\end{definition}

\begin{theorem}[The witness set contains the defect]
\label{thm:witness-contains}
Suppose exactly one edge $e^{\ast}$ is defective, with defect $D >
2\max_\ell \epsstar(\ell)$. Then for every cycle basis $\mathcal{B}$,
$e^{\ast}\in W(\mathcal{B})$, and $F(\mathcal{B})$ is precisely the set
of basis cycles containing $e^{\ast}$ an odd number of times.
\end{theorem}

\begin{proof}
A cycle's sum is $\pm D$ times its net signed traversal count of
$e^{\ast}$ and $0$ from every other edge. By \Cref{thm:detection} a
cycle is flagged iff that count is nonzero, i.e.\ iff it traverses
$e^{\ast}$ a net nonzero number of times; for simple cycles this is
containment. Hence every flagged cycle contains $e^{\ast}$, so
$e^{\ast}$ lies in their intersection, which is $W(\mathcal{B})$.
\end{proof}

\begin{corollary}[Localisation is to the witness set]
\label{cor:localise-witness}
The basis-independent content of a positive result is: the defect lies
in $W(\mathcal{B})$. The cardinality $\abs{W}$ measures the sharpness of
the localisation, with $\abs{W}=1$ the ideal and $\abs{W}$ large
indicating that the flagged cycles share many edges and the test cannot
separate them.
\end{corollary}

\begin{remark}[Minimum cycle bases sharpen the witness set]
\label{rem:mcb}
While \Cref{thm:witness-contains} holds for every basis, short cycles
intersect in fewer edges, so a minimum-weight cycle basis
\citep{horton1987,kavitha2009} typically yields a smaller $W$ and hence
a sharper localisation. We adopt one in \Cref{sec:protocol} for this
reason, and report $\abs{W}$ as a primary outcome rather than naming a
loop.
\end{remark}

% =====================================================================
\section{Validity before diagnosis}
\label{sec:validity}
% =====================================================================

A diagnostic test is meaningless on a substrate that fails its own
preconditions. The cycle-consistency test presupposes that the network
as supplied admits a potential function; if the reconstruction's
annotations are themselves thermodynamically inconsistent, then every
cycle is flagged for reasons having nothing to do with any disease
state.

\begin{definition}[Validity check]
\label{def:validity}
A network with edge data $d$ and tolerances $\epsstar$ \emph{passes the
validity check} if
$\mathrm{verdict}(\ell)=\textsf{CONSISTENT}$ for every $\ell$ in the
chosen basis, evaluated on the \emph{reference} (nominally healthy)
parameter set.
\end{definition}

\begin{criterion}[Invalid versus negative]
\label{crit:invalid}
An experiment in which the reference network fails
\Cref{def:validity} is \textbf{invalid}, not negative. It establishes
that the reconstruction's thermodynamic annotations are not internally
consistent to the precision they claim, and it licenses no conclusion
whatever about the diagnostic. An experiment in which the reference
network passes and a perturbed network is nonetheless not flagged is
\textbf{negative}, and it does license a conclusion: the test does not
detect that perturbation at that data quality.
\end{criterion}

\begin{remark}[Why this distinction must be pre-registered]
\label{rem:preregister}
Without \Cref{crit:invalid} stated in advance, a failed reference check
invites the interpretation most favourable to whoever is reporting: it
can be described as a data problem when the diagnostic succeeded and as
a theory problem when it did not. Fixing the interpretation before the
data are seen removes that degree of freedom, which is the same
methodological concern that motivates pre-registration generally
\citep{simmons2011,sandve2013}.
\end{remark}

% =====================================================================
\section{A language in which the tolerance cannot be omitted}
\label{sec:language}
% =====================================================================

The analysis above is only as good as its implementation. A
practitioner who computes $\epsstar$ correctly and then writes
\verb|if abs(H) > 1e-6| has gained nothing. We therefore express the
protocol in a small typed language in which the omission is not
expressible.

\subsection{Types}

\begin{definition}[Types]
\label{def:types}
The types are
\[
  \tau \;::=\; \cf{Circuit} \mid \cf{Loop} \mid \cf{Holonomy}
  \mid \cf{Tolerance} \mid \cf{Verdict} \mid \cf{Witness}.
\]
A value of type \cf{Holonomy} is a pair $(\widehat{\hol}_\ell,\ell)$; a
value of type \cf{Tolerance} is a triple
$(\epsnum,\epsdata,\epsstar)$ with $\epsdata$ possibly the distinguished
value \cf{undefined}; a value of type \cf{Verdict} is an element of
$\{\textsf{CONSISTENT},\textsf{UNDECIDABLE},\textsf{INCONSISTENT}\}$
\emph{paired with} the \cf{Tolerance} that produced it.
\end{definition}

\subsection{The typing rule that does the work}

\begin{rulename}[Verdict requires a tolerance]
\label{rule:verdict}
\[
  \frac{\Gamma\vdash h : \cf{Holonomy}
        \qquad \Gamma\vdash t : \cf{Tolerance}}
       {\Gamma\vdash \cf{admit}\ h\ \cf{tolerance}\ t : \cf{Verdict}}
\]
There is no other rule with \cf{Verdict} in its conclusion.
\end{rulename}

\begin{theorem}[No bare verdict]
\label{thm:no-bare-verdict}
No well-typed program produces a value of type \cf{Verdict} that is not
paired with the \cf{Tolerance} used to derive it. In particular, a
comparison of a \cf{Holonomy} against a numeric literal is not
well-typed.
\end{theorem}

\begin{proof}
By \Cref{rule:verdict} the only introduction form for \cf{Verdict}
requires a premise of type \cf{Tolerance}, and the conclusion carries it
by \Cref{def:types}. A numeric literal has no type in the grammar of
\Cref{def:types}, so there is no derivation in which one appears in
the second premise; and no coercion rule from a numeric literal to
\cf{Tolerance} exists, since \cf{Tolerance} is introduced only by the
computation of \eqref{eq:epsstar} from a \cf{Loop} and its data. Hence
no derivation produces a bare verdict.
\end{proof}

\begin{remark}[What this does and does not buy]
\label{rem:type-limits}
\Cref{thm:no-bare-verdict} rules out a class of implementation error:
it is not possible to forget the tolerance, to hard-code it, or to
report a verdict whose provenance is not recoverable. It does not make
the tolerance \emph{correct}: a program supplying a wrong $\sigma_i$ to
\eqref{eq:data-floor} is well-typed and wrong. The type system enforces
accountability of the computation, not accuracy of the inputs, and we
state this because the distinction is easy to overstate in the other
direction.
\end{remark}

\subsection{The protocol as a program}

\begin{lstlisting}[caption={The cycle-consistency test with cycle-local
tolerances. Every verdict carries its tolerance by
\Cref{thm:no-bare-verdict}; the validity gate of \Cref{crit:invalid}
precedes any diagnostic claim.},label={lst:program}]
-- kcl-tolerance.cfc
-- Cycle-consistency with per-cycle tolerances.

import Recon3D
import Recon3D.uncertainty as sigma

circuit Reference from Recon3D.glycolysis_tca {
  solve yield C_ref
}

-- Gauge-centre potentials to minimise the numerical floor
-- (Proposition 3.5). Changes no cycle sum.
let C_ref_centred := centre_potentials(C_ref)

-- Minimum-weight cycle basis (Remark 6.8)
let B := minimum_cycle_basis(C_ref_centred)

-- VALIDITY GATE (Criterion 7.2). The reference network must
-- pass before any diagnostic verdict is meaningful.
foreach loop in B {
  holonomy of loop in C_ref_centred yield h_ref
  tolerance of loop with sigma yield t
  admit h_ref tolerance t yield v_ref

  assert v_ref == CONSISTENT
    otherwise invalid
      emit "reference network fails its own consistency check"
}

emit "validity gate passed: reference is self-consistent"

-- DIAGNOSTIC. Perturbed network.
circuit Perturbed from Recon3D.glycolysis_tca {
  reaction PK : PEP -> Pyruvate, k : 0.15 * Recon3D.k_PK
  solve yield C_pert
}

let C_pert_centred := centre_potentials(C_pert)

let flagged := {}
foreach loop in B {
  holonomy of loop in C_pert_centred yield h
  tolerance of loop with sigma yield t
  admit h tolerance t yield v

  report loop, h, t.numerical, t.data, t.star, v

  where v == INCONSISTENT collect loop into flagged
}

-- Localisation is to the witness set, not to a named loop
-- (Corollary 6.7).
witness of flagged yield W
report "witness_set", W
report "witness_cardinality", size(W)

-- The claim under test, stated basis-independently.
assert Recon3D.edge("PK") in W
  emit "defect localised: PK edge in witness set"
  otherwise decline
    emit "defect not localised at this data quality"
\end{lstlisting}

% =====================================================================
\section{Protocol}
\label{sec:protocol}
% =====================================================================

\subsection{Data}

\begin{enumerate}[leftmargin=1.5em,itemsep=2pt]
\item \textbf{Network.} The glycolysis, pentose phosphate, and TCA
subnetwork of a human metabolic reconstruction
\citep{brunk2018}, approximately $80$ reactions over $60$ metabolites.
\item \textbf{Thermodynamics.} Standard transformed Gibbs energies with
reported uncertainties from a component-contribution estimator
\citep{noor2013,flamholz2012}. Species without a reported uncertainty
are recorded as such and trigger \Cref{rem:no-sigma}.
\item \textbf{Reference physiology.} Erythrocyte metabolite
concentrations and glycolytic flux from established kinetic models
\citep{rapoport1975,mulquiney1999}.
\item \textbf{Perturbation.} Pyruvate kinase activity reduced to
$15\%$ of reference, within the documented clinical range for the
deficiency \citep{zanella2005}.
\end{enumerate}

\subsection{Procedure}

\begin{protocol}[Cycle-consistency with cycle-local tolerance]
\label{prot:main}
\leavevmode
\begin{enumerate}[leftmargin=1.5em,itemsep=3pt]
\item Construct the reference circuit by
\Cref{def:mu,def:conductance}; verify the current law at every internal
node to the numerical floor of \Cref{thm:noise-scale} applied to the
node balance.
\item Gauge-centre the potentials (\Cref{prop:gauge}) and record the
resulting $\Lambda_\ell$ per cycle, both before and after.
\item Compute a minimum-weight cycle basis \citep{horton1987}; record
$\nu$, the length distribution, and the basis itself so that the
computation is reproducible.
\item For each cycle compute $\epsnum(\ell)$ by \eqref{eq:noise-bound}
and $\epsdata(\ell)$ by \eqref{eq:data-floor} at $\alpha=0.01$; record
both, and $\epsstar$.
\item \textbf{Validity gate.} Evaluate \Cref{def:validity} on the
reference. If it fails, halt and report an \emph{invalid} experiment per
\Cref{crit:invalid}.
\item Evaluate the perturbed circuit; record $\widehat{\hol}_\ell$ and
$\mathrm{verdict}(\ell)$ for every cycle.
\item Compute the witness set $W$ (\Cref{def:witness}) and record
$\abs{W}$ and whether the perturbed edge lies in it.
\item \textbf{Control 1 (null).} Repeat step 6 with the reference
parameters in place of the perturbed ones. Any \textsf{INCONSISTENT}
verdict here is a false positive and must not occur, by
\Cref{thm:no-false-positive}.
\item \textbf{Control 2 (basis).} Repeat steps 3--7 with a
\emph{different} cycle basis (e.g.\ a fundamental basis from a spanning
tree). Verify the network-level verdict agrees
(\Cref{thm:basis-verdict}) and record how $F$ and $W$ differ
(\Cref{prop:basis-localisation}).
\item \textbf{Control 3 (detection sweep).} Sweep the perturbation
magnitude and record the smallest $D$ at which the correct edge enters
$W$; compare against the predicted $2\epsstar$ of
\Cref{thm:detection}.
\end{enumerate}
\end{protocol}

\subsection{Pre-registered criteria}

\begin{criterion}[Success]
\label{crit:success}
\leavevmode
\begin{enumerate}[label=(\alph*),leftmargin=2em,itemsep=2pt]
\item \textbf{Validity.} The reference network passes
\Cref{def:validity} on every basis cycle.
\item \textbf{No false positives.} Control 1 produces zero
\textsf{INCONSISTENT} verdicts, as \Cref{thm:no-false-positive}
requires.
\item \textbf{Localisation.} The perturbed edge lies in $W$, and
$\abs{W}\le 3$.
\item \textbf{Basis invariance.} Control 2 yields the same network-level
verdict under both bases.
\item \textbf{Detection threshold.} Control 3 finds the empirical
detection threshold within a factor of $10$ of the predicted
$2\epsstar$.
\end{enumerate}
\end{criterion}

\begin{criterion}[Failure, and what each failure means]
\label{crit:failure}
\leavevmode
\begin{itemize}[leftmargin=1.5em,itemsep=2pt]
\item Failure of (a) renders the experiment \emph{invalid}
(\Cref{crit:invalid}): the reconstruction's thermodynamic annotations
are not self-consistent to their claimed precision. No conclusion about
the diagnostic follows.
\item Failure of (b) falsifies \Cref{thm:no-false-positive} as
implemented and indicates an error in the tolerance computation, not in
the theory, since the theorem is deterministic.
\item Failure of (c) with (a) and (b) passing is a genuine
\emph{negative}: the test does not localise this defect at this data
quality. If additionally most verdicts are \textsf{UNDECIDABLE}, the
limiting factor is $\epsdata$ and the finding is about the database
(\Cref{rem:undecidable}).
\item Failure of (d) would contradict \Cref{thm:basis-verdict} and
indicates an implementation error in the basis computation.
\item Failure of (e) indicates the bound of \Cref{thm:noise-scale} is
loose by more than an order of magnitude on real data, which is a
finding about the bound worth reporting.
\end{itemize}
\end{criterion}

\begin{remark}[Controls that can fail]
\label{rem:controls}
Controls 1--3 are chosen so that each can fail on well-formed input.
Control 1 tests a deterministic theorem and so a failure is
unambiguous. Control 2 tests an invariance that the theory asserts, so
a failure localises to the implementation. Control 3 tests a bound
rather than an identity, so a failure is informative about the bound's
tightness. A control that cannot fail measures nothing, and we have
excluded from the protocol several checks with that property --- notably
any check comparing the implementation's verdict against the same
implementation's tolerance, which would pass by construction.
\end{remark}

% =====================================================================
\section{Scope and limitations}
\label{sec:limits}
% =====================================================================

\paragraph{Ideal solutions.} \Cref{eq:mu} omits activity coefficients.
In the ionic strengths of cytosol these are not unity, and the
correction shifts $\mu_i$ by amounts comparable to the smaller
$\sigma_i$ \citep{alberty1998}. The effect is a systematic contribution
that our $\epsdata$, built from reported random uncertainties, does not
capture. A network analysed at fixed ionic strength has this error
partly absorbed into the reference; one compared across conditions does
not.

\paragraph{Steady state.} The circuit reading of
\Cref{prop:kirchhoff} holds at steady state. A network under transient
drive has no time-invariant $\mu$, and cycle sums computed from
instantaneous concentrations are not the quantity analysed here.

\paragraph{Independence of uncertainties.} \Cref{thm:data-floor} assumes
the $\sigma_i$ are independent. Component-contribution estimators
produce correlated errors, since species sharing chemical groups share
group contributions \citep{noor2013}. Correlation typically
\emph{reduces} the variance of a cycle sum, because the shared
components cancel around the loop, so our $\epsdata$ is conservative in
the direction of declaring more cycles \textsf{UNDECIDABLE}. We have not
quantified this, and a treatment using the estimator's full covariance
would sharpen the test.

\paragraph{Simple cycles.} \Cref{thm:data-floor} and
\Cref{thm:witness-contains} are stated for simple cycles, where
traversal counts are $\pm1$. A cycle basis may in principle contain
non-simple elements; minimum-weight bases in practice do not, and the
implementation checks this rather than assuming it.

\paragraph{One defect.} \Cref{thm:witness-contains} assumes a single
defective edge. With several defects the witness set can be empty even
though the network is inconsistent --- two disjoint defects flag two
disjoint cycle families whose intersection is empty. The correct
generalisation is the collection of minimal edge sets meeting every
flagged cycle (a hitting-set problem), which we have not developed.

\paragraph{The bound is not tight.} \eqref{eq:noise-bound} is a standard
worst-case forward bound. Random rounding errors accumulate as
$O(\sqrt{L})$ rather than $O(L)$ in practice \citep{higham2002}, so the
numerical floor is pessimistic by roughly $\sqrt{L}$. Since $\epsdata$
dominates by orders of magnitude in realistic settings
(\Cref{thm:trichotomy}(ii)), this pessimism rarely changes a verdict,
but it would matter for a network with exceptionally well-characterised
thermodynamics.

% =====================================================================
\section{Conclusion}
\label{sec:conclusion}
% =====================================================================

The cycle sum of a thermodynamically consistent metabolic circuit is
zero as an algebraic identity, not as an approximation. Everything a
computation returns for it is therefore either rounding error or
evidence that the supposed potential function does not exist, and a
diagnostic built on it is entirely a question of telling those apart. We
have argued that the customary device for doing so --- a fixed tolerance
described as a solver setting --- cannot work, because the rounding
error's scale is set by the cycle's length and potential range and so
varies by orders of magnitude across one network. We measured the
consequence: at a single fixed tolerance, the same network suffers a
$0.41$ false-positive rate on its long high-potential cycles while
missing $28\%$ of real defects on its short low-potential ones.

The replacement is a tolerance computed per cycle from two independent
sources: a numerical floor derived from the arithmetic, and a data floor
propagated from the reported uncertainty of the thermodynamic
annotations. Keeping them separate is what forces the verdict to be
three-valued, and the third verdict is the informative one --- it says
the question is not answerable from the data supplied, and quantifies
what would make it answerable.

Two further results bound what may be claimed. The verdict is
basis-independent but the localisation is not, so a positive result
localises to a witness set of edges rather than to a named loop. And a
reference network that fails its own consistency check yields an invalid
experiment rather than a negative one, a distinction that must be fixed
before the data are seen if it is to constrain anything.

What we have not done is run the protocol. It is specified here to the
level at which it can be executed without further design decisions, with
criteria stated in advance and controls chosen so that each can fail.
Whether the test diagnoses anything is a question for that execution,
and the most likely outcome --- given that $\epsdata$ exceeds $\epsnum$
by twelve orders of magnitude on current reconstructions --- is that
many cycles return \textsf{UNDECIDABLE}. That would be a finding about
the state of thermodynamic annotation in metabolic databases, and it
would be worth reporting plainly.

% =====================================================================
\section*{Reproducibility}
% =====================================================================

\Cref{prot:fixed-fails,prot:main} specify every parameter needed to
reproduce the numbers reported here. The synthetic demonstration of
\Cref{tab:fixed-fails} depends only on the stated network sizes,
potential ranges, and a seeded generator. The protocol of
\Cref{sec:protocol} depends on public data
\citep{brunk2018,noor2013,flamholz2012,rapoport1975,mulquiney1999} and
on standard implementations of minimum cycle basis
\citep{horton1987,hagberg2008} and floating-point summation
\citep{kahan1965}. No step requires proprietary software, wet-lab work,
or computation beyond a single workstation.

% =====================================================================
\bibliographystyle{plainnat}
\bibliography{references}

\end{document}
